% ============ Solutions and hints (Part 8: §22, §24) ============
% 原书 pp.295-296; 原书对 §22 的 Ex.22.1, 22.4、§23 的 Ex.23.1-23.3 未给解答, 编号按原书跳号.

\section*{\S\ref{sec:22}}

\noindent\textbf{22.2.}\quad Algebraic independence comes from Theorem~\hyperref[thm:16.2]{16.2}, (i).
To prove flatness, setting $I = \sum x_iC$ and using Theorem~\hyperref[thm:22.3]{3}, we reduce to
proving that $\Tor^C_1(k, A) = 0$. Since $x$ is a $C$-sequence, the Koszul complex
$L_\bullet = K_\bullet(x, C)$ constructed from $C$ and $x$ is a free resolution of the
$C$-module $k = C/I$, and $\Tor^C_1(k, A) = H_1(L_\bullet \otimes_C A)$. However,
$L_\bullet \otimes_C A$ is just the Koszul complex constructed from $A$ and $\bar{x}$, and since
$\bar{x}$ is an $A$-sequence, $H_1(L_\bullet \otimes A) = 0$.

\noindent\textbf{22.3.}\quad We need only show that $\Tor^A_1(k, M) = 0$. By Lemma~\hyperref[lem:18.2]{2} of
\S\ref{sec:18}, $\Tor^A_1(k, M) = \Tor^{A/xA}_1(k, M/xM)$, but by assumption the right-hand
side is $0$.

\section*{\S\ref{sec:24}}

\noindent\textbf{24.1.}\quad Use $0 \to I^i/I^{i+1} \to A/I^{i+1} \to A/I^i \to 0$ to deduce that
$\Ass(A/I^i) = \Ass(A/I)$ for all $i$.

\noindent\textbf{24.2.}\quad By Theorem~\hyperref[thm:24.5]{5}, it is enough to show that for a prime ideal $p$ of
$A$, $\CM(A/p)$ contains a non-empty open. Let $P$ be the inverse image of $p$ in $R$,
so that $A/p = R/P$. If $x_1, \ldots, x_r \in P$ are chosen to form a system of
parameters of $R$, then since $R_P$ is CM, they form an $R_P$-sequence. Thus passing
to a smaller neighbourhood of $P$, we can assume (i) $P$ is the unique minimal
prime divisor of $(x) = (x_1, \ldots, x_r)R$, and (ii) $\bar{x}$ is an $R$-sequence. Now
replacing $R$ by $R/(x)$, we can assume that $P$ is nilpotent; moreover, we can
take $P^i/P^{i+1}$ to be free $R/P$-modules. Now using the previous question it
follows easily that $R$ is CM implies $R/P$ is CM.

\noindent\textbf{24.3.}\quad After a preliminary reduction as in the previous question, use the
proof of Theorem~\hyperref[thm:24.6]{6}.
